\documentclass[../../../include/open-logic-section]{subfiles}

\begin{document}

\olfileid{sth}{z}{infinity-again}
\olsection{अनंतता}

(एखादा संच अस्तित्वात असेल, तर) अनंत संख्येने संच अस्तित्वात आहेत, हे दाखवण्यासाठी आपल्याकडे आधीच पुरेशी स्वयंसिद्धके आहेत. समजा एखादा संच अस्तित्वात आहे. मग \olref[sth][z][sep]{prop:emptyexists} नुसार $\emptyset$ अस्तित्वात असतो. आता कोणताही संच~$x$ दिल्यास \olref[sth][z][pairs]{prop:pairsconsequences} नुसार $x \cup \{x\}$ अस्तित्वात असतो. ही प्रक्रिया काही वेळा केल्यास पुढील संच मिळतात:
\begin{enumerate}
	\item[0.] $\emptyset$ %& $0$ & stage $0$\\
	\item[1.] $ \{\emptyset\}$ %& $1$ & stage $1$\\
	\item[2.] $\{\emptyset, \{\emptyset\}\}$ %& $2$ & stage $2$\\
	\item[3.] $\{\emptyset, \{\emptyset\}, \{\emptyset, \{\emptyset\}\}\}$ %&$3$ & stage $3$\\
	\item[4.] $\{\emptyset, \{\emptyset\}, \{\emptyset, \{\emptyset\}\}, \{\emptyset, \{\emptyset\}, \{\emptyset, \{\emptyset\}\}\}\}$% & $4$ & stage $4$
\end{enumerate}
हे प्रत्येक संच इतरांहून वेगळे आहेत, हे तपासता येते.

क्रमांकन $0$ पासून सुरू करण्यामागे काही कारणे आहेत. त्यांपैकी एक असे: ``$n$'' असे नाव दिलेल्या संचात नेमके $n$ सदस्य आहेत आणि सहजज्ञानानुसार तो $n$ व्या टप्प्यावर तयार होतो, हे तपासणे अवघड नाही.

पण यातून आपल्याला \emph{अनंत संख्येने} संच मिळतात; \emph{अनंत संच}, म्हणजे अनंत सदस्य असलेला एक संच, अस्तित्वात आहे याची खात्री मिळत नाही. हा फरक महत्त्वाचा आहे. कारण डेडेकिंड-अनंत संच मिळाल्याशिवाय डेडेकिंड बीजसंरचना बनवता येणार नाही. आणि नैसर्गिक संख्यांचा संच म्हणून वापरता यावी, यासाठी अशी बीजसंरचना आपल्याला हवी आहे. (\olref[sfr][infinite][dedekindsproof]{sec} शी तुलना करा.)

आतापर्यंत मांडलेली स्वयंसिद्धके कोणत्याही अनंत संचाचे अस्तित्व \emph{सिद्ध करत नाहीत}. म्हणून नवे स्वयंसिद्धक मांडावे लागेल:

\begin{axiom}[अनंतता]
असा एक संच $I$ अस्तित्वात असतो की $\emptyset \in I$ आणि $x \in I$ असेल तेव्हा $x \cup \{x\} \in I$.
\begin{align*}
	\exists I( & (\exists o \in I)\forall x\ x \notin o \land {}\\
	& (\forall x \in I)(\exists s \in I)\forall z(z \in s \liff (z \in x \lor z = x)))
\end{align*}
\end{axiom}

अनंततेच्या स्वयंसिद्धकातून मिळणारा संच $I$ डेडेकिंड-अनंत आहे, हे सहज दिसते. त्याचा निवडक घटक $\emptyset$ आहे आणि $I$ वरील एकास-एक फलन $s(x) = x\cup
\{x\}$ असे आहे. आता \olref[sfr][infinite][dedekind]{thm:DedekindInfiniteAlgebra} मध्ये डेडेकिंड-अनंत संचापासून डेडेकिंड बीजसंरचना कशी मिळते ते पाहिले; तिलाच आपण नैसर्गिक संख्यांचा संच मानू. अधिक नेमकेपणे:
\begin{defn}\ollabel{defnomega}
अनंततेच्या स्वयंसिद्धकातून मिळालेला कोणताही संच $I$ घ्या. $s$ हे $s(x) = x \cup \{x\}$ असे फलन असू द्या. $\omega =
\closureofunder{s}{\emptyset}$ असे मानू. $\omega$ चे सदस्य \emph{नैसर्गिक संख्या} म्हणू. $n$ ही $\emptyset$ वर $s$ च्या $n$ वेळा केलेल्या उपयोगाचे फलित आहे, असे म्हणू.
\end{defn}

आधीच्या काही परिच्छेदांत ``$n$'' अशी खूण केलेला संच आता $n$ ही संख्याच \emph{मानला} जाईल, हे मागे पाहून तपासता येते.

या निर्णयाचे महत्त्व \olref[sth][z][nat]{sec} मध्ये पाहू. सध्या त्याच्या आधारे एक सहज वाटणारा निष्कर्ष सिद्ध करता येतो:

\begin{prop}\ollabel{naturalnumbersarentinfinite}
कोणतीही नैसर्गिक संख्या डेडेकिंड-अनंत नाही.
\end{prop}

\begin{proof}
सिद्धता विगमनाने करायची आहे; \olref[sfr][infinite][induction]{thm:dedinfiniteinduction} पाहा. स्पष्टच, $0
= \emptyset$ डेडेकिंड-अनंत नाही. विगमनाच्या पायरीसाठी आपण प्रतिपक्ष सिद्ध करू: $s(n)$ डेडेकिंड-अनंत असल्याचे (अशक्य असूनही) मानल्यास $n$ डेडेकिंड-अनंत ठरेल.

म्हणून $s(n)$ डेडेकिंड-अनंत आहे असे मानू. म्हणजे असा एखादा !!{injection} $f$ आहे की $\ran{f}\subsetneq \dom{f} = s(n) = n \cup
\{n\}$. आता दोन प्रकरणे पाहावी लागतील.

\emph{प्रकरण १: $n \notin \ran{f}$.} मग $\ran{f} \subseteq n$ आणि $f(n) \in n$. $g = \funrestrictionto{f}{n}$ मानू; मग $\ran{g} =
\ran{f} \setminus \{f(n)\} \subsetneq n = \dom{g}$. म्हणून $n$ डेडेकिंड-अनंत आहे.

\emph{प्रकरण २: $n \in \ran{f}$.} एखादा $m \in \dom{f} \setminus \ran{f}$ निश्चित करू आणि प्रांत $s(n) = n \cup \{n\}$ असलेले फलन $h$ पुढीलप्रमाणे ठरवू:
\[
h(x) =
	\begin{cases}
		f(x) & \text{जर }f(x) \neq n\\
		m & \text{जर }f(x)=n
	\end{cases}
\]
म्हणून $h$ आणि $f$ सर्वत्र जुळतात; फक्त $h(f^{-1}(n)) = m
\neq n = f(f^{-1}(n))$ या ठिकाणी वेगळे आहेत. $f$ हे !!a{injection} असल्याने $n \notin
\ran{h}$; तसेच $\ran{h} \subsetneq \dom{h} = s(n)$. आता प्रकरण १ मधील युक्तिवादाने $n$ डेडेकिंड-अनंत ठरतो.
\end{proof}

पण अनंततेच्या स्वयंसिद्धकाचे \emph{समर्थन} कसे करायचे, हा प्रश्न उरतो. थोडक्यात उत्तर असे की आपण सांगत आलेल्या मांडणीत आणखी एक तत्त्व जोडावे लागेल. ते असे:
\begin{enumerate}
\item[] \stagesinf. अनंत टप्पा असतो. म्हणजे असा टप्पा असतो की (a) तो पहिला नसतो, (b) त्याच्या आधी काही टप्पे असतात आणि (c) त्याचा लगतचा पूर्ववर्ती टप्पा नसतो.
\end{enumerate}
या तत्त्वावरून अनंततेचे स्वयंसिद्धक सरळ मिळते. नैसर्गिक संख्या $n$ ही टप्पा $n$ वर तयार होते, हे आपल्याला ठाऊक आहे. म्हणून $\omega$ हा संच पहिल्या अनंत टप्प्यावर तयार होतो. $\omega$ स्वतःच अनंततेच्या स्वयंसिद्धकासाठी अपेक्षित उदाहरण ठरतो.

मात्र यामुळे \stagesinf{} चे समर्थन कसे करायचे हा प्रश्नच पुढे ढकलला जातो. \stagessucc{} प्रमाणे हे तत्त्वही संचांच्या संचयी-क्रमिक पदानुक्रमाबद्दलच्या मूळ मांडणीत स्पष्टपणे नव्हते. त्याहून महत्त्वाचे म्हणजे संच टप्प्याटप्प्याने तयार होतात या कल्पनेतूनच प्रक्रियेत \emph{अनंत} टप्पा असलाच पाहिजे, असे निष्पन्न होत नाही. टप्प्यांचा क्रम नैसर्गिक संख्यांसारखा आहे, असे मानणे पूर्णपणे सुसंगत वाटते.

मग एक उघड अडचण उभी राहते. \olref[sfr][infinite][dedekindsproof]{sec} मध्ये विचारांचा डेडेकिंड-अनंत संच आहे या डेडेकिंड यांच्या ``सिद्धतेचा'' आपण विचार केला. ती तुम्हाला फार समाधानकारक वाटली नसेल. पण संचांच्या क्रमिक संकल्पनेतून (किंवा ``विचारांच्या नियमां''तून) \stagesinf{} आपोआप स्वीकारणे भाग पडत नसेल, तर डेडेकिंड-अनंत संच अस्तित्वात आहे या दाव्याचे अंतर्गत समर्थन अजूनही मिळत नाही.

या विषयावर आणखी बरेच काही सांगता येईल. पण एखादे स्वयंसिद्धक किंवा तत्त्व समर्थित करण्यासाठी \emph{काय लागते}, याचा विचार सुरू करता येईल अशा टप्प्यावर तुम्ही आता पोहोचला असाल, अशी आशा आहे. पुढे आपण \stagesinf{} गृहीत धरू.

\end{document}
